\documentclass[aspectratio=169]{beamer}

\usetheme{Madrid}
\usecolortheme{default}
\usefonttheme{professionalfonts}

\usepackage{amsmath,amssymb,bm}
\usepackage{braket}
\usepackage{listings}
\usepackage{hyperref}

% ---------------------------
% Listings
% ---------------------------
\lstset{
  basicstyle=\ttfamily\small,
  columns=fullflexible,
  breaklines=true,
  frame=single,
  showstringspaces=false
}

\title{Exercises week 4}
\subtitle{Exercises with tentative solutions}
\author{Try to do the exercises first without looking at the solutions}
\institute{}
\date{}

\begin{document}

% ==========================================================
\begin{frame}
  \titlepage
\end{frame}


\begin{frame}{Exercise 1: Partially entangled two-qubit state}
Consider
\[
\ket{\psi(\theta)}=\cos\theta\,\ket{00}+\sin\theta\,\ket{11},\qquad \theta\in[0,\pi/2].
\]
\textbf{Tasks:}
\begin{enumerate}
\item Compute $\rho_{AB}=\ket{\psi(\theta)}\bra{\psi(\theta)}$.
\item Compute $\rho_A=\Tr_B(\rho_{AB})$.
\item Find eigenvalues of $\rho_A$ and compute
\[
S(\rho_A),\quad S_2(\rho_A),\quad S_\alpha(\rho_A).
\]
\item Identify $\theta$ for which entanglement is maximal.
\end{enumerate}

\medskip
\textbf{Hint:} $\rho_A$ will be diagonal in $\{\ket{0},\ket{1}\}$.
\end{frame}

% ==========================================================
\begin{frame}{Solution 1 (full derivation)}
Start with
\[
\rho_{AB}=\ket{\psi}\bra{\psi}=
\cos^2\theta\,\ket{00}\bra{00}
+\cos\theta\sin\theta\,\ket{00}\bra{11}
+\cos\theta\sin\theta\,\ket{11}\bra{00}
+\sin^2\theta\,\ket{11}\bra{11}.
\]

Use $\Tr_B(\ket{ab}\bra{cd})=\braket{d}{b}\ket{a}\bra{c}$:
\begin{align}
\Tr_B(\ket{00}\bra{00})&=\ket{0}\bra{0},\\
\Tr_B(\ket{00}\bra{11})&=\braket{1}{0}\ket{0}\bra{1}=0,\\
\Tr_B(\ket{11}\bra{00})&=\braket{0}{1}\ket{1}\bra{0}=0,\\
\Tr_B(\ket{11}\bra{11})&=\ket{1}\bra{1}.
\end{align}

Therefore,
\[
\rho_A=\cos^2\theta\,\ket{0}\bra{0}+\sin^2\theta\,\ket{1}\bra{1}
=\begin{pmatrix}
\cos^2\theta & 0\\
0 & \sin^2\theta
\end{pmatrix}.
\]

Eigenvalues: $\lambda_1=\cos^2\theta$, $\lambda_2=\sin^2\theta$.
\end{frame}

% ==========================================================
\begin{frame}{Solution 1: Entropies}
Von Neumann entropy:
\[
S(\rho_A)=-\cos^2\theta\,\log_2(\cos^2\theta)-\sin^2\theta\,\log_2(\sin^2\theta).
\]

Second Rényi entropy:
\[
S_2(\rho_A)=-\log_2\Tr(\rho_A^2)
=-\log_2\left(\cos^4\theta+\sin^4\theta\right).
\]

General Rényi:
\[
S_\alpha(\rho_A)=\frac{1}{1-\alpha}\log_2\left((\cos^2\theta)^\alpha+(\sin^2\theta)^\alpha\right).
\]

\medskip
\textbf{Max entanglement:} at $\theta=\pi/4$,
$\lambda_1=\lambda_2=1/2$ so $S(\rho_A)=1$ bit.
\end{frame}

% ==========================================================
\begin{frame}{Exercise 2: Measurement entropy vs state entropy}
Let the single-qubit state be
\[
\rho = p\ket{0}\bra{0}+(1-p)\ket{1}\bra{1},\qquad p\in[0,1].
\]

\textbf{Tasks:}
\begin{enumerate}
\item Compute $S(\rho)$.
\item Compute measurement entropy $H_{\text{meas}}$ for measuring in the $Z$ basis.
\item Compute $H_{\text{meas}}$ for measuring in the $X$ basis.
\end{enumerate}

\medskip
\textbf{Key point:} $S(\rho)$ is basis-independent, while $H_{\text{meas}}$ depends on the measurement.
\end{frame}

% ==========================================================
\begin{frame}{Solution 2: $Z$ and $X$ measurements}
Since $\rho$ is diagonal with eigenvalues $(p,1-p)$:
\[
S(\rho)=-p\log_2 p-(1-p)\log_2(1-p).
\]

\medskip
\textbf{$Z$ basis measurement:}
\[
p(0)=p,\quad p(1)=1-p
\Rightarrow
H_Z=-p\log_2 p-(1-p)\log_2(1-p)=S(\rho).
\]

\medskip
\textbf{$X$ basis measurement:}
$\ket{\pm}=(\ket{0}\pm\ket{1})/\sqrt{2}$, with projectors $\Pi_\pm=\ket{\pm}\bra{\pm}$.
Then
\[
p(\pm)=\Tr(\rho\,\Pi_\pm)=\frac{1}{2},
\]
so
\[
H_X=1,
\]
independent of $p$.

\medskip
Thus generally $H_{\text{meas}}\ge S(\rho)$, with equality in eigenbasis of $\rho$.
\end{frame}

% ==========================================================
\begin{frame}{Exercise 3: Mutual information for a classical mixture}
Consider a two-qubit mixed state
\[
\rho_{AB}= \frac{1}{2}\ket{00}\bra{00}+\frac{1}{2}\ket{11}\bra{11}.
\]

\textbf{Tasks:}
\begin{enumerate}
\item Compute $\rho_A,\rho_B$.
\item Compute $S(\rho_{AB})$, $S(\rho_A)$, $S(\rho_B)$.
\item Compute $I(A\!:\!B)=S(\rho_A)+S(\rho_B)-S(\rho_{AB})$.
\item Compare to a Bell state, where $I(A\!:\!B)=2$.
\end{enumerate}
\end{frame}

% ==========================================================
\begin{frame}{Solution 3: Classical vs quantum correlations}
Compute reduced states:
\[
\rho_A=\rho_B=\frac{I_2}{2}\quad\Rightarrow\quad S(\rho_A)=S(\rho_B)=1.
\]

The global state has eigenvalues $\{1/2,1/2,0,0\}$, hence
\[
S(\rho_{AB})=1.
\]

Therefore
\[
I(A\!:\!B)=1+1-1=1.
\]

\medskip
\textbf{Interpretation:}
\begin{itemize}
\item This state has correlations, but they are purely classical.
\item Bell states have maximal mutual information $I=2$, reflecting strong quantum correlations.
\end{itemize}
\end{frame}

\end{document}
